\begin{derivation}[从\eqnrefs{eq:sm-full-lagrangian-blocks-dp11}到\eqnrefs{eq:minimal-sm-parameter-count-dp11}]
Yukawa 矩阵的矩阵元依赖味基选择，参数计数必须先除去这种场重定义冗余，因此从夸克扇区开始。三代时，$Y_u$ 与 $Y_d$ 各是一般复 $3\times3$ 矩阵。每个矩阵含

\begin{equation*}
 9\ \text{个复矩阵元}
 =18\ \text{个实参数},
\end{equation*}
所以两者合计有 $36$ 个实参数。

暂时令 Yukawa 耦合为零。夸克动能项在世代空间具有
\begin{equation*}
 U(3)_{Q_L}\times U(3)_{u_R}\times U(3)_{d_R}
\end{equation*}
的味基冗余。每个 $U(3)$ 有 $3^2=9$ 个连续生成元，因此三组基变换一共有 $27$ 个连续参数。不过不能直接用 $36-27$，因为一般非零的 $Y_u,Y_d$ 仍保留共同的矢量式重子数变换
\begin{equation*}
 Q_L\to \ee^{\ii\alpha}Q_L,
 \qquad
 u_R\to \ee^{\ii\alpha}u_R,
 \qquad
 d_R\to \ee^{\ii\alpha}d_R.
\end{equation*}
在 Yukawa 双线性项中，左手场经 Dirac 共轭后产生 $\ee^{-\ii\alpha}$，恰与右手场的 $\ee^{+\ii\alpha}$ 相消。因此 $U(1)_B$ 是打开 Yukawa 耦合后仍未破缺的子群，不能用来改变 $Y_u,Y_d$。真正能够消去坐标冗余的味基参数只有
\begin{equation*}
 27-1=26.
\end{equation*}
所以夸克 Yukawa 扇区留下
\begin{equation*}
 36-26=10
\end{equation*}
个物理参数。为了看清这 $10$ 个参数为什么恰好表现为“六个质量、三个角、一个相位”，还要把上述商空间计数落实到质量本征基。对两个一般复矩阵分别作奇异值分解：
\begin{equation*}
 Y_u=U_{uL}\,y_u\,U_{uR}^\dagger,
 \qquad
 Y_d=U_{dL}\,y_d\,U_{dR}^\dagger,
\end{equation*}
其中
\begin{equation*}
 y_u=\operatorname{diag}(y_u,y_c,y_t),
 \qquad
 y_d=\operatorname{diag}(y_d,y_s,y_b)
\end{equation*}
可取为非负实对角矩阵。右手幺正矩阵 $U_{uR},U_{dR}$ 把两个质量矩阵对角化；带电弱流只含左手双重态，因此两个左手旋转之间的相对矩阵保留下来：
\begin{equation*}
 V_{\rm CKM}=U_{uL}^\dagger U_{dL}.
\end{equation*}
六个奇异值与六个夸克质量一一对应，剩下的问题就是一个一般 $3\times3$ 幺正矩阵究竟含多少可观测参数。

一般 $U(3)$ 有 $9$ 个实参数。可把它们组织成三个实混合角与六个相位。质量项已经对角以后，六个 Dirac 夸克仍可分别作共同左右手重相位
\begin{equation*}
 u_i\to \ee^{\ii\alpha_i}u_i,
 \qquad
 d_j\to \ee^{\ii\beta_j}d_j,
\end{equation*}
而不改变任何质量本征值。带电流
$\bar u_{Li}\gamma^\mu V_{ij}d_{Lj}$ 因而只把 CKM 元素变换为
\begin{equation*}
 V_{ij}\longrightarrow
 \ee^{-\ii\alpha_i}V_{ij}\ee^{+\ii\beta_j}.
\end{equation*}
六个相位 $\{\alpha_i,\beta_j\}$ 中，若全部同时增加同一个 $\omega$，则
\begin{equation*}
 \ee^{-\ii(\alpha_i+\omega)}V_{ij}
 \ee^{+\ii(\beta_j+\omega)}
 =\ee^{-\ii\alpha_i}V_{ij}\ee^{+\ii\beta_j},
\end{equation*}
所以共同的重子数相位并不能改变 $V_{\rm CKM}$。真正可用于消去 CKM 相位的只有 $6-1=5$ 个独立重相位。因此六个原始相位中只剩
\begin{equation*}
 6-5=1
\end{equation*}
个物理 $CP$ 相位，而三个实混合角保留。于是夸克 Yukawa 扇区的 $10$ 个参数被显式分解为
\begin{equation*}
 \boxed{6\ \text{个夸克质量}
 +3\ \text{个 CKM 混合角}
 +1\ \text{个 CKM 相位}.}
\end{equation*}

带电轻子扇区只有一个一般复 $3\times3$ 矩阵 $Y_e$，起始共有 $18$ 个实参数。令 Yukawa 耦合为零时，动能项具有
\begin{equation*}
 U(3)_{L_L}\times U(3)_{e_R}
\end{equation*}
的 $18$ 参数味基冗余。严格的最小标准模型中没有右手中微子，中微子保持无质量；把 $Y_e$ 对角化以后，三个轻子家族仍可分别做保持动能项和质量项不变的共同矢量式重相位。因此未破缺子群是 $U(1)^3$，维数为 $3$。可用于消去 $Y_e$ 坐标冗余的基变换参数数目为
\begin{equation*}
 18-3=15,
\end{equation*}
于是只剩
\begin{equation*}
 18-15=3
\end{equation*}
个带电轻子质量；在这个严格最小场内容中没有物理的 PMNS 混合角与相位。

所以全部 Yukawa 与味结构一共贡献
\begin{equation*}
 10+3=13
\end{equation*}
个物理参数。再加入三种规范耦合 $g_s,g,g'$、Higgs 势中的两个参数 $\mu_H^2,\lambda_H$，以及一个基不变的强相互作用 $CP$ 参数 $\bar\theta$，总数为
\begin{equation*}
 13+3+2+1=19.
\end{equation*}
于是得到\eqnrefs{eq:minimal-sm-parameter-count-dp11}。19 个连续参数就是耦合参数空间对场重定义轨道取商后的维数；粒子质量、混合参数和规范耦合只是这个商空间的一组物理坐标。
\end{derivation}
